%\DocumentMetadata{xmp=false,uncompress}%xmp=true,uncompress,
%\DocumentMetadata{pdfversion=2.0,_pdfstandard/UA-2,_pdfstandard/A-4,_pdfstandard/X-6n,tagging=on,lang=en-US}%xmp=true,
\DocumentMetadata{pdfversion=2.0,pdfstandard={a-4,ua-2},tagging=on,lang=en-US}%xmp=true,
%\\DocumentMetadata{pdfversion=1.7,pdfstandard=x-1b,lang=en-US}%xmp=true,ua-1
%\RequirePackage[l2tabu, orthodox]{nag} not with 2.0

\documentclass[a4paper, USenglish]{article}%scrartcl [12pt, a4paper]%, english]


%\usepackage[rerun=always]{pythontex}%rerun=always

%\usepackage[destlabel,pdfusetitle]{hyperref}
%\usepackage{listings}

% TBD: eliminate later 
\PassOptionsToPackage{pdfusetitle}{hyperref}

\input{../header.tex}

\ifPDFTeX%
\else
\usepackage{unicode-math} % REIE: for lualatex or xelatex only 
\fi
% TBD: work out 


\IfPackageLoadedTF{tex4ht}{
}{
  \newindex[LaTeX Packages]{pkg}
  \newindex[LaTeX Commands]{cmd}
  \makeindex
}

\IfPackageLoadedTF{tex4ht}{
  \newcommand{\cmd}[2][]{\texttt{\textbackslash{}#2#1}}
  \newcommand{\pkg}[1]{\texttt{#1}}% without indexing 
  % \newcommand{\tool}[1]{\texttt{#1}}% without indexing 
  % \newcommand{\param}[1]{\texttt{#1}}% without indexing 
}
{
  %TBD: ask in formum: would prefer texorpdfstring in the macto. 
  %\newcommand{\cmd}[2][]{\texorpdfstring{\texttt{\textbackslash{}#2#1}\sindex[cmd]{#2@\texttt{\textbackslash{}#2}}}{xxx}}%\unichar{0092}
 \newcommand{\cmd}[2][]{\texttt{\textbackslash{}#2#1}\sindex[cmd]{#2@\texttt{\textbackslash{}#2}}}%\unichar{0092}
%   \newcommand{\cmd}[2][]{%
%   \ifpdfstring
%     \unichar{92}#2%
%   \else
%     \texttt{\textbackslash{}#2#1}\sindex[cmd]{#2@\texttt{\textbackslash{}#2}}%
%   \fi
% }
  % TBD: clarify different behavior: 1 or two arguments. 
  % seemingly there is a problem with detokenize also 
  \newcommand{\pkg}[2][]{\texorpdfstring{\texttt{#2}\sindex[pkg]{#2@\texttt{#2}}}{xxx}} % TBD: this must be extracted 
  % \newcommand{\tool}[1]{\texttt{#1}\sindex[tool]{#1}}% TBD: this must be extracted 
  % \newcommand{\param}[1]{\texttt{#1}\sindex[param]{#1}}% TBD: this must be extracted 
}

%\usepackage{hyperxmp}

\input{../headerSuppressMetaPDF.tex}

\hypersetup{%
  pdfinfo={
    %Author      ={Ernst Reissner},
    %Title       ={On PDF files created from LaTeX and analyzing them },
    %Creator     ={unknown}, 
    % % LaTeX with hyperref; this is uniform for hyperref. 
    % % except for beamer class, then it is LaTeX with Beamer class
    % % else TeX for lua/pdflatex, but  XeTeX output 2023.12.06:0106 for xelatex 
    % % thus it is advisable to eliminate this 
    % % maybe security issue, but it is well visible. 
    % % also corrupts reproducibility. 
    % Producer    ={unknown},
    % % xdvipdfmx with version and that like for dvi or xelatex; 
    % % pdfTeX-1.40.25, LuaTeX-1.17.0 for pdf
    % % This shall be eliminated. 
    %CreationDate={unknown},
    % ModDate     ={unknown},
    % % these two are only overwritten to avoid wrong entry 
    % % when using \texttt{SOURCE\_DATE\_EPOCH=0 FORCE\_SOURCE\_DATE=1}
    % % This is recommended for tests, because else we cannot give a \date 
    Subject     ={Form and analysis of PDF files created from LaTeX},
    Keywords    ={LaTeX;PDF;pdfinfo;exiftool;pdflatex;xelatex;lualatex}
  }
  %pdfsource   ={no latex main file}
}

%\newcommand{\cmd}[1]{\textbackslash #1}

\title{On PDF files created from \LaTeX{} and analyzing them  }
\author{Ernst Reissner (rei3ner@arcor.de)}
\date{\Today}

\begin{document}

\maketitle
\tableofcontents

% TBD: glossary based on a database. 
% It is quite natural that this 

\section{Introduction}\label{sec:intro}

This document is created with \texttt{lualatex} or that like 
with output format 
\ifpdf%
PDF%
\else
DVI% TBD: take XDV into account also 
\fi.
The package \texttt{tex4ht} 
is \IfPackageLoadedTF{tex4ht}{}{not} loaded. 


This is about the special form of PDF files created from \LaTeX{} 
and about analyzing them. 
Accordingly, there are two sections,~\ref{sec:pdfFromLatex} on the special form of PDF files 
if created from \LaTeX{} and~\ref{sec:pdfAna} on how to analyze the created PDF file. 
Note that this document covers both, PDF 1.7 described in~\cite{Pdf17} as a representative of PDF 1.x 
and PDF 2.0 described in~\cite{Pdf20}. 

This document is a torso, Section~\ref{sec:pdfAna} contains a lot of experimental material. 
Maybe, we shall reorganize, merging creation and analysis and treating various aspects one after the other. 


\section{PDF files from \LaTeX}\label{sec:pdfFromLatex}

Modern \LaTeX{} compilers are mainly focussed on creating PDF files. 
On the other hand, PDF technology is quite general 
whereas \LaTeX{} compilers are just one creator for PDF files out of many. 
PDF technology evolves quickly and the \LaTeX{} world struggles to keep pace. 
So in general, \LaTeX{} compilers do not implement all features offered by PDF technology. 
Furthermore, \LaTeX{} is very special, and thus the created files will be very special forever. 

Whereas the visual appearance of PDF file from \LaTeX{} sources is quite mature, 
the focus is now on metadata which is data on the content, so to speak describing the nature of the various pieces of information 
independent of their representation. 

\subsection{Metadata comprising XMP}\label{subsec:meta}

% TBD: this requires rework. 
% First of all, we shall explain what XMP data is and that latex creates it if \DocumentMetadata is configured alike. 
% This could be a separate section. 

% After that maybe a separate section on time stamps with the pieces of information below. 

\subsection{Metadata with \texorpdfstring{\cmd{DocumentMetadata}}{\unichar{0092}DocumentMetadata}}\label{subsec:docMeta}
%\subsection{Metadata with \texorpdfstring{\cmd{DocumentMetadata}}{\unichar{0092}DocumentMetadata}}\label{subsec:docMeta}

This precedes usage of package \pkg{hyperref} as described in Section~\ref{subsec:hyperref} because takes over certain functionality: 
The options \texttt{pdfa}, \texttt{pdfversion} should no longer be used. 
Others interfere like \texttt{CJKbookmarks}, \texttt{psdextra} and \texttt{nesting}. 
Take care of \texttt{booktabs}, which may change. 
For a collection of all settings and starting point for further pieces of information see~ \cite{HyperTextP}, Section 5.10. 

Interesting is option \texttt{pdfusetitle} which ``tries to'' take over title and also author. 





\subsection{The package \texorpdfstring{\pkg{hyperref}}{hyperref}}\label{subsec:hyperref}

Here we describe how to set author, title, subject and keywords using \texttt{hyperref}. 

% Another aspect deserving analysis is the features of package hyperref 
% and how we use it to include pieces of information but also how to hide other aspects for security reasons. 
% provided: Title, author, subject, keywords 
% hidden: creator, producer 
% possibly forced: creation date and related. 

One aspect is that \LaTeX{} creates PDF from TEX independently of old versions of the PDF file. 
Thus, it does not use the intrinsic version management system. 
As a consequence, (initial) creation date and modification date, i.e.\@ date of last modification coincide. 
The same for last modification date of metadata if XMP data is written. 
%  <xmp:CreateDate>2025-12-14T12:59:53Z</xmp:CreateDate>
%  <xmp:ModifyDate>2025-12-14T12:59:53Z</xmp:ModifyDate>
%  <xmp:MetadataDate>2025-12-14T12:59:53Z</xmp:MetadataDate>
%  CreationDate:    Sun Dec 14 13:59:53 2025 CET
%  ModDate:         Sun Dec 14 13:59:53 2025 CET


\subsection{Hiding metadata with injection \texttt{headerSuppressMetaPDF.tex}}\label{subsec:hide}



\section{Analysis of PDF files}\label{sec:pdfAna}

Two tools for analyzing general metadata in PDF files deserve to be highlighted, \texttt{exiftool} and \texttt{pdfinfo}. 
Currently, we restrict ourselves to the latter and describe it in Section~\ref{subsec:pdfAnaPdfInfo}. 
In addition, \texttt{show-pdf-tags} displays the tags in a tagged PDF file. 
This is described in Section~\ref{subsec:pdfAnaShowTags}. 


\subsection{The \texttt{pdfinfo} tool}\label{subsec:pdfAnaPdfInfo}

Let us start with metadata of a TEX file without \cmd{DocumentMetadata}, loaded package \pkg{hyperref}. 
Without any options, \texttt{pdfinfo} yields something like the following: 

\begin{verbatim}
Creator:         TeX
Producer:        LuaTeX-1.22.0
CreationDate:    Sun Dec 14 20:47:26 2025 CET
ModDate:         Sun Dec 14 20:47:26 2025 CET
Custom Metadata: yes
Metadata Stream: no
Tagged:          no
UserProperties:  no
Suspects:        no
Form:            none
JavaScript:      no
Pages:           3
Encrypted:       no
Page size:       595.276 x 841.89 pts (A4)
Page rot:        0
File size:       45572 bytes
Optimized:       no
PDF version:     1.7
\end{verbatim}

TBD\@: clarify what we see here really. 
Is it a particular dictionary? 

Page size is not specific for PDF created from TEX, but it is due to the option \texttt{a4paper} 
passed to \cmd{documentclass}; else this would be \texttt{letter}. 

Note that \texttt{Creator} and \texttt{Producer} is specific for PDF created from TEX or even for the \LaTeX{} compiler. 
The \texttt{xelatex} compiler deviates from the others insofar as the creator is xelatex, but the producer is the backend as e.g. \texttt{xdvipdfmx}. 

Also, that \texttt{CreationDate} and \texttt{ModDate} coincide is typical for a tool chain based on \LaTeX, 
because \LaTeX{} compiler are not incremental: they do not add diffs to original PDF files, 
but create anew. 

What is also typical are that in the basic configuration, most values are \texttt{no} or \texttt{none}. 
In particular, there is no \texttt{Metadata Stream} and consequently, invoking \texttt{pdfinfo} with option \texttt{-meta} is empty. 
Likewise, the document is not \texttt{Tagged} and consequently, invoking \texttt{show-pdf-tags} yields 
%
\begin{verbatim}
Document catalog has no markinfo.Marked entry. It might not be tagged.
\end{verbatim}
%
Still the return value is $0$. 

Tagging information can be retrieved by \texttt{pdfinfo -struct} and by \texttt{pdfinfo -struct-text}, 
but we decided to treat these functions together with \texttt{show-pdf-tags} in Section~\ref{subsec:pdfAnaShowTags}. 
TBD here> combination with other options, \texttt{opw}, \texttt{upw}



Only \texttt{Custom Metadata} is set. 
As a consequence, invoking \texttt{pdfinfo} with option \texttt{-custom} is not empty. 
In our case it shows 
%
\begin{verbatim}
CreationDate:    Sun Dec 14 21:08:46 2025 CET
Creator:         TeX
ModDate:         Sun Dec 14 21:08:46 2025 CET
PTEX.FullBanner: This is LuaHBTeX, Version 1.22.0 (TeX Live 2025)
Producer:        LuaTeX-1.22.0
\end{verbatim}
%
but the only novelty is the banner. 
It is even specific for the compiler. 
%PTEX.Fullbanner: This is pdfTeX, Version 3.141592653-2.6-1.40.28 (TeX Live 2025) kpathsea version 6.4.1
Note that \texttt{xelatex} creates much different output and in particular no banner. 

By the way, the date format can be modified with options \texttt{-isodates} and \texttt{-rawdates}. 
Both options can be well combined with option \texttt{-custom}. 
Note that the help text on \texttt{-custom} is misleading. 

Loading the package \pkg{hyperref} 
adds the keys \texttt{Title} \texttt{Subject} \texttt{Keywords}  and \texttt{Author} to the metadata, 
but the according values are empty. 
Loading with option \texttt{pdftitle}, seems to include the title and the author given by \cmd{title} and \cmd{author}, respectively. 
Note that the \texttt{Creator} is no longer just \texttt{TeX}, but \texttt{aTeX with hyperref}. 

Also, custom metadata becomes richer, namely Author, Keywords, Subject and Title added: 
%
\begin{verbatim}
Author:          Ernst Reissner (rei3ner@arcor.de)
CreationDate:    Sun Dec 14 22:09:44 2025 CET
Creator:         LaTeX with hyperref
Keywords:        LaTeX;PDF;pdfinfo;exiftool;pdflatex;xelatex;lualatex
ModDate:         Sun Dec 14 22:09:44 2025 CET
PTEX.FullBanner: This is LuaHBTeX, Version 1.22.0 (TeX Live 2025)
Producer:        LuaTeX-1.22.0
Subject:         Form and analysis of PDF files created from LaTeX
Title:           On PDF files created from LaTeX and analyzing them 
\end{verbatim}
%
Interesting enough, setting these user properties, 
does keeps the value for key \texttt{UserProperties} unchanged at value \texttt{no}. 
% TBD: search in hyperref documentation: there is a way to add further properties. 
% Maybe this changes the value. 

Note that including \texttt{headerSuppressMetaPDF.tex} suppresses the custom metadata and more than that. 
As a consequence, key \texttt{Custom Metadata} receives value \texttt{no} and consequently 
\texttt{pdfinfo} invoked with option \texttt{-custom} is empty. 
This hides hints on the toolchain which may be a security measure. 

Now we specify 
%
\begin{verbatim}
\DocumentMetadata{xmp=false,pdfversion=2.0,lang=en-US}
\end{verbatim}
%

The default is \texttt{xmp=true}, so if removing, the value for key \texttt{Metadata Stream} switches to \texttt{yes} 
and accordingly, \texttt{pdfinfo} invoked with option \texttt{-meta} is no longer empty. 

TBD: describe here. Maybe better extract. It has a lot of details. 

Accordingly, if setting \texttt{tagging} to \texttt{on} activates tagging of course 
and so the key \texttt{Tagged} receives the value \texttt{yes}. 
If tagging is not ok, then the suspect flag may be set, i.e.\@ the key \texttt{Suspects} receives the value \texttt{yes}. 
TBD\@: more research based on the specifications~\cite{Pdf17} and~\cite{Pdf20}. 
It is also not clear whether this can be set, although certain PDF standards are satisfied, 
like that on Well-Tagged-PDF~\cite{WtPDF1}. 


Note that package \pkg{hyperref} supports creating forms also. 
For details see~\cite{HyperTextP}, Section 9. % TBD: eforms package
Also, there is the \texttt{eforms} package described in~\cite{eformP}. 
TBD\@: It must be checked whether then the key \texttt{Form} receives the value \texttt{yes}. 

% \documentclass{article}
% \usepackage[pdftex]{insdljs}
 
% \OpenAction{\JS{%
% app.alert("HALLO WELT");
% }}
 
% \begin{document}
% THIS DOCUMENT IS A DOCUMENT
% \end{document}



% Creator:         TeX
% Producer:        LuaTeX-1.22.0
% CreationDate:    2025-12-14T19:47:26Z
% ModDate:         2025-12-14T19:47:26Z
% Custom Metadata: yes
% Metadata Stream: no
% Tagged:          no
% UserProperties:  no
% Suspects:        no
% Form:            none
% JavaScript:      no
% Pages:           3
% Encrypted:       no
% Page size:       595.276 x 841.89 pts (A4)
% Page rot:        0
% File size:       45572 bytes
% Optimized:       no


Applying \texttt{pdfinfo} with option \texttt{-isodates} on the PDF file created from this TEX file yields 
something like the following. \\

\noindent
\begin{verbatim}
Title: On PDF files created from and analyzing them
Subject: Form and analysis of PDF files created from LaTeX
Keywords: LaTeX;PDF;pdfinfo;exiftool;pdflatex;xelatex;lualatex
Author: Ernst Reissner (rei3ner@arcor.de)
Creator: unknown
Producer: unknown
CreationDate: 2025-12-25T21:15:21Z
ModDate: 2025-12-25T21:15:21Z
Custom Metadata: no
Metadata Stream: yes
Tagged: yes
UserProperties: no
Suspects: no
Form: none
JavaScript: no
Pages: 8
Encrypted: no
Page size: 595.276 x 841.89 pts (A4)
Page rot: 0
File size: 133298 bytes
Optimized: no
PDF version: 1.7
\end{verbatim}
% \begin{pycode}
% import subprocess

% output = subprocess.check_output("pdfinfo -isodates pdfFromLatex.pdf", shell=True)
% str = output.decode("utf-8")
% str = str.replace("\n", "\\\\")
% print(str)
% \end{pycode}

As one can see, the lines take the form \texttt{<key>: <value}, 
but the original is even better: The space after the colon is in fact a tab 
so that the values are printed in line. 

Output of that kind is created without any option, with options referring to passwords like \texttt{-opw} and \texttt{-opw}. 
Alone or in conjunction with options referring to date like \texttt{-isodates} or \texttt{-rawdates}, 
the format of the dates can be adapted to the own needs. 

Before moving on to the further options, 
let us describe the meaning of the keys in detail: 
%
\begin{description}
\item[Title] The title set by package \texttt{hyperref} as described in Section~\ref{subsec:hyperref}
\item[Subject] The subject set by package \texttt{hyperref} as described in Section~\ref{subsec:hyperref}
\item[Keywords] The keywords set by package \texttt{hyperref} as described in Section~\ref{subsec:hyperref}
\item[Author] The author(s) set by package \texttt{hyperref} as described in Section~\ref{subsec:hyperref}
\end{description}





\subsection{The \texttt{show-pdf-tags} tool}\label{subsec:pdfAnaShowTags}

Here, it must be clarified the relation between \texttt{show-pdf-tags}, \texttt{pdfinfo -struct} and \texttt{pdfinfo -struct-text}

\subsection{The \texttt{qpdf} tool}\label{subsec:pdfAnaQpdf}

This provides a low level detailed sight onto the PDF file. 
%
\begin{Verbatim}
qpdf --show-object=trailer pdfFromLatex.pdf 
\end{Verbatim}
%
shows the trailer directory as described in~/cite{Pdf20}, Table~15. 
Inserting line breaks this looks as follows: 
%
\begin{Verbatim}
<< 
/Filter /FlateDecode 
/ID [ <2192177301c209a1f9a0efaf865496a0> <2192177301c209a1f9a0efaf865496a0> ] 
/Index [ 0 574 ] 
/Info 572 0 R 
/Length 1401 
/Root 571 0 R 
/Size 574 
/Type /XRef 
/W [ 1 3 1 ] 
>>
\end{Verbatim}

Strictly speaking, the trailer IDs are optional in many situations, but in fact, without, the PDF document is quite degraded. 
According to~\cite{Pdf20}, Section~14.4, it consists of two hashes. 
The two coincide if the document is newly created, 
whereas only the second one is adapted with incremental updates. 
Since \LaTeX{} compilers do not support internal versioning, the two hashes always coincide. 
It is quite unfortunate that the specification how to create the hashes is a ``may'' condition. 
It may depend on 
%
\begin{itemize}
\item the current time, 
\item a string representation of file location, 
\item the size of the file. 
\end{itemize}

The author hopes that the \LaTeX{} compiler conform with this and even that it is not the file location 
but only the file name. 
There is a lack of documentation. 
For \texttt{pdflatex},~\cite{PdfTexUsr25}, 
Section~4.2.9 indicates that by default only the processing time and the file name are used. 



The only aspect we are currently interested in is the trailer ID which consists of two hashes. 
PDF has an 

\bibliographystyle{alpha}
\bibliography{../lit,../litTools}{}% chktex 11 

\IfPackageLoadedTF{tex4ht}{%
  % no index and no glossaries with tex4ht 
}{
%\chapter{Indices}
  %\printindex[idx]
  %\printindex[tool]
  \printindex[pkg]
  \printindex[cmd]
  %\printindex[param]
  %\printglossaries%
}

\end{document}
